\documentclass[11pt]{article}
\usepackage[margin=1in]{geometry}
\usepackage{graphicx}
\usepackage{amsmath}

\title{CSCI 311: Connecting a City\\
Minimum Spanning Tree Algorithm Comparison}
\author{CityConnections Project}
\date{}

\begin{document}

\maketitle

\section{Introduction}

This project addresses the problem of connecting all intersections in a city using the minimum total length of cable. The road network is modeled as a weighted graph where nodes represent intersections and edges represent roads with distances. The goal is to compute a Minimum Spanning Tree (MST), which connects all nodes with the smallest possible total edge weight.

In addition to correctness, this project evaluates the performance of multiple MST algorithms to understand trade-offs in runtime and implementation design.

\section{Algorithms}

We implemented three different MST algorithms:

\subsection{Heap-Based Prim’s Algorithm}
This version of Prim’s algorithm uses a priority queue (min-heap) to efficiently select the next minimum edge connecting the current tree to a new node. The algorithm maintains a set of visited nodes and repeatedly adds the smallest edge that expands the tree.

\textbf{Time Complexity:} $O(E \log V)$

\subsection{Naive Prim’s Algorithm}
This version uses a simpler approach by scanning all edges that connect the current tree to the rest of the graph. It uses vectorized operations (NumPy) instead of explicit nested loops.

\textbf{Time Complexity:} $O(V^2)$

\subsection{Kruskal’s Algorithm}
Kruskal’s algorithm sorts all edges by weight and adds them one by one, ensuring no cycles are formed. A Disjoint Set Union (DSU) structure is used to efficiently detect cycles.

\textbf{Time Complexity:} $O(E \log E)$

\section{Implementation Details}

The graph is stored using a pandas DataFrame containing edge information. For Heap Prim’s, an adjacency list is constructed to efficiently access neighboring edges. A custom heap implementation is used to maintain priority ordering.

Kruskal’s algorithm uses a Disjoint Set Union data structure with path compression and union-by-rank to optimize performance.

The Naive Prim’s algorithm uses NumPy operations to efficiently filter and select edges without explicit loops.

\section{Results}

Each algorithm was run on the same dataset, and runtime was measured using a high-resolution timer. All three algorithms produced MSTs with the same number of edges ($V - 1$) and identical total weight, confirming correctness.

The observed runtime behavior matched expectations:

\begin{itemize}
\item Heap Prim’s performed the fastest due to efficient edge selection using a heap.
\item Kruskal’s performed well but incurred overhead from sorting edges.
\item Naive Prim’s was significantly slower due to repeated scanning of edges.
\end{itemize}

\section{Discussion}

The results demonstrate that algorithmic design significantly impacts performance. While all three approaches solve the same problem correctly, their efficiency differs due to underlying data structures and operations.

Heap Prim’s achieves better performance by reducing unnecessary work through priority queue operations. Kruskal’s benefits from global sorting but can be slower on dense graphs. The Naive Prim’s approach is simpler but less efficient, especially as graph size increases.

These findings align with theoretical time complexities, highlighting the importance of choosing appropriate algorithms and data structures for large-scale problems.

\section{Conclusion}

This project illustrates how multiple algorithms can solve the same problem with varying efficiency. By implementing and comparing three MST algorithms, we gained insight into both theoretical and practical performance differences. Efficient data structures such as heaps and disjoint sets play a critical role in improving runtime, particularly for large graphs.

\end{document}
